\section{Въведение. Видове сходимости на случайни величини}

Законите за големите числа са едни от най-фундаменталните резултати в теорията на вероятностите. Често пъти, когато използваме различни методи, компютърни пакети или статистически анализи (като например анализ на данни), обосновката за валидността на тези методи произтича именно от този фундамент. За тяхното формулиране и разбиране първо е необходимо понятието за сходимост в теорията на вероятностите.

Ще разгледаме три основни типа сходимост на случайни величини. Съществуват и други видове, но тези три са най-важните за нашите цели. За първите два типа сходимост ще допуснем, че имаме редица от случайни величини $X_n$ (за $n \ge 1$) и една гранична случайна величина $X$. Важно е да отбележим, че за тези първи две дефиниции всички случайни величини живеят в едно и също вероятностно пространство $(\Omega, \mathcal{A}, \mathbb{P})$. Тоест, имаме едно общо множество от елементарни събития $\Omega$, една сигма-алгебра $\mathcal{A}$ и една вероятностна мярка $\mathbb{P}$.

\subsection{Сходимост почти сигурно (п.с.)}

За удобство, нека означим с $L_X$ множеството от всички елементарни събития $\omega \in \Omega$, за които числовата редица $X_n(\omega)$ се схожда към числото $X(\omega)$ в класическия смисъл на математическия анализ:
\[
L_X = \left\{ \omega \in \Omega \mathrel{\bigg|} \lim_{n \to \infty} X_n(\omega) = X(\omega) \right\}
\]
За всяко конкретно $\omega$, $X_n(\omega)$ е просто една редица от реални числа. Ако тази числова редица клони към $X(\omega)$, то събитието $\omega$ принадлежи на $L_X$.

\begin{defn}[Сходимост почти сигурно]
Казваме, че редицата $X_n$ се схожда почти сигурно към $X$, което се записва като $X_n \xrightarrow{\text{п.с.}} X$, ако вероятността на събитието $L_X$ е равна на единица:
\[
\mathbb{P}(L_X) = \mathbb{P}\left( \lim_{n \to \infty} X_n = X \right) = 1
\]
\end{defn}

Интуитивно това означава, че ако изберете случайно едно елементарно събитие (т.е. проведете експеримент), вероятността да изберете такова $\omega$, че редицата да се сходи, е $1$. Това не означава задължително, че сходимостта е налице за абсолютно всяко възможно $\omega$ — множеството от изключения (където няма сходимост) може да не е празно, но то задължително има вероятност (мярка) 0.

\subsection{Сходимост по вероятност}

За втория тип сходимост нека въведем събитието $A_{n, \varepsilon}$, което се състои от всички елементарни събития, за които разликата между $X_n$ и $X$ надвишава дадено число $\varepsilon > 0$:
\[
A_{n, \varepsilon} = \{ \omega \in \Omega \mathrel{\big|} |X_n(\omega) - X(\omega)| > \varepsilon \}
\]
Това са онези изходи от експеримента, при които стойността на $X_n$ стои на разстояние по-голямо от $\varepsilon$ спрямо стойността на $X$.

\begin{defn}[Сходимост по вероятност]
Казваме, че редицата $X_n$ се схожда към $X$ по вероятност (означава се с $X_n \xrightarrow{\mathbb{P}} X$), ако за всяко $\varepsilon > 0$ е изпълнено:
\[
\lim_{n \to \infty} \mathbb{P}(|X_n - X| > \varepsilon) = 0
\]
\end{defn}

Това означава, че при достатъчно голямо $n$, вероятността разстоянието между $X_n$ и $X$ да бъде съществено (повече от $\varepsilon$) става пренебрежимо малка. Вероятността не е задължително да намалява монотонно, но в граница тя клони към нула.

\emph{Забележка:} В горната дефиниция е напълно достатъчно да се разглеждат само $\varepsilon$ от вида $\frac{1}{r}$ за всяко естествено число $r \ge 1$. Ако $\varepsilon$ е произволно реално число, винаги можем да намерим естествено $r$, такова че $\frac{1}{r+1} < \varepsilon \le \frac{1}{r}$. Тъй като условието $|X_n - X| > \varepsilon$ влече $|X_n - X| > \frac{1}{r+1}$, събитията се влагат едно в друго (т.е. $A_{n, \varepsilon} \subseteq A_{n, \frac{1}{r+1}}$). Работата с изброимо количество стойности на $\varepsilon$ е много полезна в теоретичните доказателства.

\subsection{Сходимост по разпределение}

При предишните два вида сходимост беше ключово всички случайни величини да са дефинирани в едно и също вероятностно пространство $\Omega$. Сходимостта по разпределение обаче е най-слабият, но може би най-полезният тип сходимост в практиката, тъй като не зависи от конкретното пространство. Тя се интересува единствено от вероятностните свойства (разпределенията) на величините.

Нека $F_X(x)$ е функцията на разпределение на случайната величина $X$. Ще означим с $C_{F_X}$ множеството от всички реални числа $x \in \mathbb{R}$, в които функцията $F_X$ е непрекъсната. За непрекъснати случайни величини (с плътност), $C_{F_X} = \mathbb{R}$. Ако обаче $X$ има скокове (дискретна случайна величина, напр. на Бернули), $C_{F_X}$ е цялата реална права без точките на скок. Точките на скок са точно тези $x$, за които $\mathbb{P}(X = x) > 0$.

\begin{defn}[Сходимост по разпределение]
Нека $X_n$ е случайна величина в пространство $(\Omega_n, \mathcal{A}_n, \mathbb{P}_n)$ за $n \ge 1$, а $X$ е случайна величина в пространство $(\Omega, \mathcal{A}, \mathbb{P})$. Казваме, че $X_n$ се схожда по разпределение към $X$ (записваме $X_n \xrightarrow{d} X$), ако за всяко $x \in C_{F_X}$ е вярно, че:
\[
\lim_{n \to \infty} F_{X_n}(x) = F_X(x)
\]
\end{defn}

Този вид сходимост означава, че въпреки че експериментите може да са коренно различни (едното изследване може да е върху хора, друго върху животни, трето хвърляне на монети), вероятностната структура на случайните величини става все по-близка. Това е сходимост на самите функции на разпределение, а не сходимост на реализациите на величините.

\section{Връзка между видовете сходимост}

Нека сега градираме видовете сходимост по тяхната „сила“.

\begin{thm}
\begin{enumerate}
    \item[а)] Ако $X_n \xrightarrow{\text{п.с.}} X$, то $X_n \xrightarrow{\mathbb{P}} X$.
    \item[б)] Ако $X_n \xrightarrow{\mathbb{P}} X$, то $X_n \xrightarrow{d} X$.
\end{enumerate}
\end{thm}

Най-силната сходимост е „почти сигурно“, следвана от „по вероятност“, и най-слабата е „по разпределение“. Поради това сходимостта по разпределение често се постига най-лесно.

Обратните твърдения по принцип не са верни, и двата контрапримера си заслужават да се видят.

\begin{example}[По вероятност, но не почти сигурно]\label{ex:10-typewriter}
Работим върху $\Omega = [0,1]$ с $\mathbb{P}$ равна на дължината. Разглеждаме индикаторите на
последователни интервали, които обикалят отрезъка:
\[
\xi_{n,k} = \ind_{\left[\frac{k-1}{n},\, \frac{k}{n}\right]},
\qquad k = 1, \dots, n, \quad n = 1, 2, \dots
\]
и ги подреждаме в една редица: $\xi_{1,1};\ \xi_{2,1}, \xi_{2,2};\ \xi_{3,1}, \xi_{3,2}, \xi_{3,3};\ \dots$

Тази редица клони към $0$ \emph{по вероятност}: за фиксирано $\varepsilon \in (0,1)$
събитието $\{\xi_{n,k} > \varepsilon\}$ е точно интервалът с дължина $\frac1n$, тоест
неговата вероятност е $\frac1n \to 0$.

Но никъде не клони \emph{почти сигурно}: за всяко фиксирано $\omega \in [0,1]$ и за всяко
$n$ съществува точно едно $k$, за което $\omega$ попада в съответния интервал, тоест
$\xi_{n,k}(\omega) = 1$. Значи редицата приема стойност $1$ безкрайно често и стойност $0$
безкрайно често — тя не се схожда за нито едно $\omega$.
\end{example}

\begin{example}[По разпределение, но не по вероятност]\label{ex:10-samedist}
Нека $\xi, \xi_1, \xi_2, \dots$ са независими и еднакво разпределени, всички $N(0,1)$.
Понеже всички имат една и съща функция на разпределение, тривиално
$\xi_n \xrightarrow{d} \xi$.

Сходимост по вероятност обаче няма. Разликата $\xi_n - \xi$ е сума на две независими
стандартно нормални величини (с обратен знак), тоест $\xi_n - \xi \sim N(0,2)$ за всяко $n$.
Тогава $\mathbb{P}(|\xi_n - \xi| > \varepsilon)$ е една и съща положителна константа за всяко
$n$ и не клони към нула.

Този пример показва и защо сходимостта по разпределение не се интересува от
вероятностното пространство: тя гледа само разпределенията, не и самите реализации.
\end{example}

\begin{proof}[Скица на доказателството на част а)]
Знаем, че $\mathbb{P}(L_X) = 1$, където $L_X = \left\{ \lim_{n\to\infty} X_n = X \right\}$. По дефиниция за сходимост на числова редица, множеството $L_X$ може да бъде представено чрез сечения и обединения:
\[
L_X = \bigcap_{r=1}^{\infty} \bigcup_{n=1}^{\infty} \bigcap_{k=n}^{\infty} A_{k, \frac{1}{r}}^{c}
\]
където $A_{k, \frac{1}{r}}^{c} = \left\{ |X_k - X| \le \frac{1}{r} \right\}$. Изразът отразява факта, че „за всяко $r$ (аналог на $\varepsilon$), съществува $n$, такова че за всяко $k \ge n$, отклонението $|X_k - X|$ е не повече от $\frac{1}{r}$“.

Взимайки допълнението по правилата на Де Морган, получаваме множеството, в което няма сходимост:
\[
L_X^c = \bigcup_{r=1}^{\infty} \underbrace{\bigcap_{n=1}^{\infty} \bigcup_{k=n}^{\infty} A_{k, \frac{1}{r}}}_{B_r}
\]
Тъй като $\mathbb{P}(L_X) = 1$, то $\mathbb{P}(L_X^c) = 0$. От това следва, че:
\[
0 = \mathbb{P}\left(\bigcup_{r=1}^{\infty} B_r\right) \ge \mathbb{P}(B_\ell) \quad \forall \ell \ge 1
\]
Следователно $\mathbb{P}(B_r) = 0$ за всяко $r \ge 1$.

Нека дефинираме $C_{n, r} = \bigcup_{k=n}^{\infty} A_{k, \frac{1}{r}}$. Забелязваме, че с нарастването на $n$, обединението включва все по-малко множества, така че редицата е намаляваща: $C_{n, r} \supseteq C_{n+1, r} \supseteq \dots$
Тъй като $B_r = \bigcap_{n=1}^{\infty} C_{n, r}$ и редицата е монотонно намаляваща, от непрекъснатостта на вероятностната мярка следва:
\[
\mathbb{P}(B_r) = 0 \implies \lim_{n\to\infty} \mathbb{P}(C_{n, r}) = 0
\]
Но $C_{n, r}$ очевидно съдържа първия си елемент $A_{n, \frac{1}{r}} = \left\{ |X_n - X| > \frac{1}{r} \right\}$. Щом вероятността на по-голямото множество $C_{n, r}$ клони към 0, тогава и:
\[
\lim_{n\to\infty} \mathbb{P}\left(|X_n - X| > \frac{1}{r}\right) = 0
\]
Това е вярно за всяко $r$, което доказва, че $X_n$ клони към $X$ по вероятност.

\end{proof}

\begin{proof}[Скица на доказателството на част б)]
Тук искаме да свържем събитието $\{X_n < x\}$, което участва във функцията на разпределение
на $X_n$, със събитието $\{X < x\}$. Директна връзка няма — трябва да минем през малка
пертурбация $\varepsilon$. Означаваме
\[
A_{n,\varepsilon} = \{ |X_n - X| > \varepsilon \},
\qquad\text{като}\qquad
\mathbb{P}(A_{n,\varepsilon}) \xrightarrow[n \to \infty]{} 0
\]
за всяко $\varepsilon > 0$ по условие (това е точно сходимостта по вероятност).

Разбиваме $\{X_n < x\}$ по $A_{n,\varepsilon}$ и допълнението му, което не променя нищо:
\[
\{X_n < x\} = \big(\{X_n < x\} \cap A_{n,\varepsilon}^c\big) \cup \big(\{X_n < x\} \cap A_{n,\varepsilon}\big).
\]
Върху $A_{n,\varepsilon}^c$ имаме $|X_n - X| \le \varepsilon$, тоест $X$ се различава от $X_n$ с
най-много $\varepsilon$. Затова
\[
\{X_n < x\} \subseteq \{X < x + \varepsilon\} \cup A_{n,\varepsilon},
\qquad
\{X < x - \varepsilon\} \cap A_{n,\varepsilon}^c \subseteq \{X_n < x\}.
\]
Прилагайки вероятност към двете влагания, получаваме „сандвич“:
\[
F_X(x-\varepsilon) - \mathbb{P}(A_{n,\varepsilon})
\ \le\ F_{X_n}(x) \ \le\
F_X(x+\varepsilon) + \mathbb{P}(A_{n,\varepsilon}).
\]
При $n \to \infty$ членовете $\mathbb{P}(A_{n,\varepsilon})$ отпадат и остава
\[
F_X(x-\varepsilon) \ \le\ \liminf_n F_{X_n}(x) \ \le\ \limsup_n F_{X_n}(x) \ \le\ F_X(x+\varepsilon).
\]
Накрая пускаме $\varepsilon \to 0$. Ако $x$ е точка на непрекъснатост на $F_X$, двата края
клонят към $F_X(x)$ и следователно $F_{X_n}(x) \to F_X(x)$ — точно това е сходимост по
разпределение.

\end{proof}

\subsection{Обратна връзка в частни случаи}
В общия случай от сходимост по разпределение не следва сходимост по вероятност. Съществува обаче едно важно изключение: когато граничната случайна величина е константа.

\begin{prop}
Ако $X_n \xrightarrow{d} c$, където $c$ е константа, то $X_n \xrightarrow{\mathbb{P}} c$ (разбира се, допускайки, че случайните величини са дефинирани в едно и също вероятностно пространство).
\end{prop}

\emph{Идея на доказателството:} При използваната тук конвенция функцията на разпределение на константата $c$ е $F_c(x) = 0$ за $x \le c$ и $F_c(x) = 1$ за $x > c$. Единствената точка на прекъсване е $x=c$. Тъй като $X_n \xrightarrow{d} c$, за всяко $x \neq c$ имаме сходимост на функциите на разпределение. От това може да се покаже, че вероятността $X_n$ да се отклони от $c$ с повече от $\varepsilon$ се състои от сумата на вероятностите $\mathbb{P}(X_n > c + \varepsilon)$ и $\mathbb{P}(X_n < c - \varepsilon)$. И двете клонят към 0, тъй като в точките $c+\varepsilon$ и $c-\varepsilon$ функцията на разпределение $F_c$ е съответно 1 и 0, към които се схождат $F_{X_n}$.

\section{Неравенство на Чебишов}

Когато въведохме дисперсията, обсъдихме, че тя играе ролята на средноквадратична грешка. Неравенството на Чебишов отговаря на въпроса: „Ако знаем нещо за дисперсията, каква е горната граница на вероятността разликата между $X$ и очакването на $X$ да е по-голяма от някакво число $a$?“ Това ни позволява лесно да оценим вероятността за големи отклонения.

\begin{prop}[Неравенство на Марков]\label{prop:markov}
Ако $Y \ge 0$ е неотрицателна случайна величина с крайно очакване, то
\[ \mathbb{P}(Y > a) \le \frac{\E Y}{a} \quad \forall a > 0 . \]
\end{prop}

\begin{proof}
Понеже $Y \ge 0$, за всяко $a>0$ е в сила поточковото неравенство
$Y \ge a\,\ind_{\{Y > a\}}$: върху събитието $\{Y>a\}$ лявата страна е по-голяма от $a$, а
извън него дясната страна е нула. Взимаме очакване и използваме монотонността му заедно с
$\E \ind_{\{Y>a\}} = \mathbb{P}(Y>a)$:
\[ \E Y \ge a\,\mathbb{P}(Y > a) . \]
\end{proof}

Неравенството на Чебишов се получава от него с едно заместване: прилагаме
твърдение~\ref{prop:markov} към $Y = (X - \E X)^2 \ge 0$ и към прага $a^2$. По-долу е
дадено и прякото доказателство, както беше направено на лекцията.

\begin{keythm}[Неравенство на Чебишов]
Нека $X$ е случайна величина с дисперсия $\Var X$. Тогава:
\[
\mathbb{P}(|X - \E X| > a) \le \frac{\Var X}{a^2} \quad \forall a > 0
\]
\end{keythm}

\begin{proof}
Нека положим $Y = X - \E X$. Търсим да оценим $\mathbb{P}(|Y| > a)$, което може да се запише като очакване на индикаторната функция: $\E[\ind_{\{|Y| > a\}}]$.
От свойствата на дисперсията имаме $\Var X = \Var Y = \E[Y^2]$ (тъй като $\E Y = 0$).

Представяме единицата като сума от два индикатора: $1 = \ind_{\{|Y| \le a\}} + \ind_{\{|Y| > a\}}$. Тогава:
\[
\E[Y^2] = \E[Y^2 \ind_{\{|Y| \le a\}}] + \E[Y^2 \ind_{\{|Y| > a\}}]
\]
Тъй като и двете събираеми са неотрицателни, можем да премахнем първото и да получим неравенство:
\[
\E[Y^2] \ge \E[Y^2 \ind_{\{|Y| > a\}}]
\]
Върху събитието, за което индикаторът $\ind_{\{|Y| > a\}}$ не е нула (тоест когато $|Y| > a$), със сигурност е изпълнено $Y^2 > a^2$. Затова можем да заменим $Y^2$ с $a^2$:
\[
\E[Y^2 \ind_{\{|Y| > a\}}] \ge \E[a^2 \ind_{\{|Y| > a\}}] = a^2 \E[\ind_{\{|Y| > a\}}] = a^2 \mathbb{P}(|Y| > a)
\]
Окончателно получаваме $\Var X \ge a^2 \mathbb{P}(|Y| > a)$, или:
\[
\mathbb{P}(|X - \E X| > a) \le \frac{\Var X}{a^2}
\]
\end{proof}

\subsection{Следствия и пример}
Аналогично неравенство за произволни моменти $m \ge 1$ се получава, като приложим твърдение~\ref{prop:markov} към $Y = |X - \E X|^m$ и прага $a^m$ (названието Марков или Чебишов тук не се фиксира, тъй като няма съществена разлика в типа доказателство):
\[
\mathbb{P}(|X - \E X| > a) \le \frac{\E[|X - \E X|^m]}{a^m}
\]
Когато $m=2$, получаваме класическото неравенство на Чебишов.

\begin{example}\label{ex:10-1}
Каква е вероятността една случайна величина $X$ да се отклони от очакването си с повече от $10$ пъти своето стандартно отклонение?
Нека за удобство $\E X = 0$. Търсим $\mathbb{P}(|X| > 10\sqrt{\Var X})$. В неравенството на Чебишов избираме $a = 10\sqrt{\Var X}$. Получаваме:
\[
\mathbb{P}(|X| > 10\sqrt{\Var X}) \le \frac{\Var X}{(10\sqrt{\Var X})^2} = \frac{\Var X}{100 \Var X} = \frac{1}{100}
\]
Това е изключително силен универсален резултат — за \emph{всяка} случайна величина вероятността да надвиши 10 пъти стандартното си отклонение никога не надхвърля $1\%$.

\end{example}

\section{Закони за големите числа (ЗГЧ)}

Законите за големите числа формализират интуитивната идея, че при извършване на голям брой независими експерименти, средноаритметичното от резултатите се приближава до теоретичното математическо очакване.

\subsection{Слаб закон за големите числа}

\begin{keythm}[Слаб ЗГЧ]
Нека $(X_i)_{i=1}^\infty$ е редица от независими и еднакво разпределени случайни величини, такива че $\E|X_1| < \infty$. Тогава за редицата е в сила законът за големите числа, тоест:
\[
\frac{\sum_{i=1}^n X_i}{n} \xrightarrow[n \to \infty]{\mathbb{P}} \E X_1
\]
\emph{Бележка:} Тъй като $X_i$ са еднакво разпределени, всички те имат едно и също математическо очакване, $\E X_i = \E X_1$.
\end{keythm}

\begin{proof}[Доказателство (в случая, когато дисперсията съществува)]
Ще докажем теоремата при допълнителното по-силно условие, че $\Var X_1 < \infty$, тъй като тогава можем директно да използваме неравенството на Чебишов.
Трябва да покажем, че за всяко $\varepsilon > 0$:
\[
\lim_{n \to \infty} \mathbb{P} \left( \left| \frac{\sum_{i=1}^n X_i}{n} - \E X_1 \right| > \varepsilon \right) = 0
\]
Тъй като $\E X_i = \E X_1$ за всяко $i$, можем да внесем очакването вътре в сумата:
\[
\mathbb{P} \left( \left| \frac{\sum_{i=1}^n X_i}{n} - \frac{\sum_{i=1}^n \E X_i}{n} \right| > \varepsilon \right) = \mathbb{P} \left( \frac{\left| \sum_{i=1}^n (X_i - \E X_i) \right|}{n} > \varepsilon \right)
\]
Като умножим двете страни на неравенството в скобите по $n$, получаваме:
\[
\mathbb{P} \left( \left| \sum_{i=1}^n (X_i - \E X_i) \right| > \varepsilon n \right)
\]
Сега прилагаме неравенството на Чебишов за случайната величина $Y = \sum_{i=1}^n (X_i - \E X_i)$, чието очакване е 0. Като заместим $a = \varepsilon n$, имаме:
\[
\mathbb{P}(|Y| > \varepsilon n) \le \frac{\Var \left( \sum_{i=1}^n (X_i - \E X_i) \right)}{(\varepsilon n)^2}
\]
Тук ключово използваме \emph{независимостта} на случайните величини $X_i$. Поради нея, дисперсията на сумата е равна на сумата от дисперсиите:
\[
\frac{\sum_{i=1}^n \Var(X_i)}{\varepsilon^2 n^2}
\]
Но тъй като величините са еднакво разпределени, те имат еднаква дисперсия $\Var X_1$. Сумата се състои от $n$ такива събираеми:
\[
= \frac{n \Var X_1}{n^2 \varepsilon^2} = \frac{\Var X_1}{n \varepsilon^2}
\]
Тъй като дисперсията $\Var X_1$ и $\varepsilon^2$ са фиксирани константи, при $n \to \infty$ този израз клони към $0$. С това доказахме сходимостта по вероятност.
\end{proof}

\subsection{Усилен закон за големите числа (УЗГЧ)}

Слабият закон гарантира сходимост само по вероятност. Усиленият закон утвърждава много по-силен резултат — сходимост почти сигурно.

\begin{keythm}[Усилен закон за големите числа]\label{thm:slln}
Нека $(X_i)_{i=1}^\infty$ е редица от независими и еднакво разпределени случайни величини с крайни очаквания $\E|X_1| < \infty$. Тогава средното аритметично се схожда почти сигурно към теоретичното очакване:
\[
\frac{\sum_{i=1}^n X_i}{n} \xrightarrow[n \to \infty]{\text{п.с.}} \E X_1
\]
Това е най-полезният и често срещан вид на закона, тъй като в повечето практически задачи (като например събиране на емпирични данни) работим с независими наблюдения от един и същи феномен.
\end{keythm}

\subsection{Примери и приложения на ЗГЧ}

\begin{example}[Схема на Бернули: гласуване или ваксини]\label{ex:10-2}
Нека $X_i \sim \Ber(p)$ с $p \in (0,1)$. Например, $X_i = 1$, ако $i$-тият човек гласува за партия А (или ако ваксината проработи при него) и $0$ в противен случай. Не знаем истинската вероятност $p$, но чрез УЗГЧ знаем, че:
\[
\frac{\sum_{i=1}^n X_i}{n} \xrightarrow[n \to \infty]{\text{п.с.}} p = \E X_1
\]
Тоест, ако попитаме достатъчно голям брой хора, пропорцията от гласове (средноаритметичното на $X_i$) ще клони почти сигурно към истинската вероятност $p$.
\end{example}

\begin{example}[Проверка на грамаж]\label{ex:10-3}
Ако на етикета пише, че в един сандвич има 50 грама шунка, можем да измерим грамажа на голям брой $n$ сандвичи. Нека $X_i$ е количеството шунка в $i$-тия сандвич. Средното аритметично $\frac{1}{n}\sum X_i$ трябва да клони към теоретичното очакване $\E X_1$. Ако измерим средно 46 грама при достатъчно голямо $n$, по силата на ЗГЧ имаме сериозни основания да се усъмним, че истинското очакване е 50 грама.
\end{example}

\begin{example}[Игра на рулетка]\label{ex:10-4}
Връщаме се към \crosslecture{ex:05-3}{примера с рулетката от лекция 5}. Печелите 1 лев с вероятност $\frac{18}{37}$ и губите 1 лев с вероятност $\frac{19}{37}$. Очакваната ви печалба от една игра е:
\[
\E X_1 = 1 \cdot \frac{18}{37} - 1 \cdot \frac{19}{37} = -\frac{1}{37}
\]
Колкото повече играете, толкова средната ви печалба на игра клони към:
\[
\frac{\sum_{i=1}^n X_i}{n} \xrightarrow[n \to \infty]{\text{п.с.}} -\frac{1}{37}
\]
което обяснява защо дългосрочно винаги сте на загуба.
\end{example}

\begin{example}[Дърпане на въже: случайно блуждаене]\label{ex:10-5}\label{ex:tug-of-war}
Връщаме се към \crosslecture{ex:05-4}{модела със случайно дърпане на въже от лекция 5}. Имаме $n$ играчи (или молекули), които дърпат въжето наляво или надясно с вероятност $\frac{1}{2}$. Силата $X_i = 1$ за надясно и $-1$ за наляво. Очакваната сила е $\E X_i = 0$. Резултантната сила е $S_n = \sum_{i=1}^n X_i$. Според ЗГЧ, средната сила $\frac{S_n}{n}$ клони към 0. Важно е обаче да не се бърка сходимостта на средната сила с вероятността дясната страна да победи ($\mathbb{P}(S_n > 0)$). Пропорцията $\frac{S_n}{n}$ клони към 0, но вероятността резултантната сила да е положителна клони към $\frac{1}{2}$, а не към 0! Това е често допускана грешка от неразбиране на разликата между самите пропорции и вероятностите за техните стойности.
\end{example}

\begin{example}[Симулации Монте Карло]\label{ex:10-6}\label{ex:monte-carlo}
Продължаваме \crosslecture{ex:monte-carlo-volume}{геометричната интерпретация на Монте Карло от лекция 2}. Искаме да намерим обема $|A|$ на дадено множество $A$, вложено в единичния куб. Алгоритъмът Монте Карло „хвърля“ случайно и равномерно точки $X_i$ в единичния куб. Дефинираме $Y_i = 1$, ако $X_i \in A$, и $Y_i = 0$, ако $X_i \notin A$. Величината $Y_i$ е Бернулиева с параметър $|A|$ (тъй като обемът на единичния куб е 1). Според УЗГЧ:
\[
\frac{\sum_{i=1}^n Y_i}{n} \xrightarrow[n \to \infty]{\text{п.с.}} |A|
\]
Тоест, просто като броим точките, попаднали в $A$, и ги разделим на общия брой хвърлени точки, ние числено апроксимираме търсения обем $|A|$. Колко точно точки трябва да хвърлим, за да постигнем определена точност, е въпрос, на който отговаря Централната Гранична Теорема, която ще разгледаме в следващите лекции.
\end{example}

\section{Задачи}

\begin{enumerate}
    \item Докажете, че в дефиницията за сходимост по вероятност е достатъчно условието да се проверява за $\varepsilon=1/r$, $r\in\mathbb{N}$.
    \item Довършете доказателството, че от $X_n\xrightarrow{\mathbb{P}}X$ следва $X_n\xrightarrow{d}X$.
    \item Довършете доказателството на частния обратен резултат: ако $X_n\xrightarrow{d}c$ за константа $c$, то $X_n\xrightarrow{\mathbb{P}}c$.
    \item Нека $\E X=0$ и $\Var X<\infty$. Използвайте неравенството на Чебишов, за да докажете универсалната оценка
    \[
    \mathbb{P}\left(|X|>10\sqrt{\Var X}\right)\le\frac{1}{100}.
    \]
\end{enumerate}
